\chapter{Introduction}\label{sec:introduction}
The Phillips curve, originally proposed by \textcite{phillips1958} and later 
formalised by \textcite{samuelson1960}, suggests an inverse relationship 
between inflation and unemployment. Despite its theoretical appeal, the 
empirical stability of this relationship has been widely debated. The curve 
appeared to break down during the stagflation of the 1970s, and more recently, 
the post-COVID surge in inflation has reignited the discussion of whether 
unemployment retains meaningful predictive power over inflation dynamics 
\parencite{stock2019}. Understanding whether labour market conditions carry 
forecasting value for inflation is not merely an academic question - it has 
direct implications for monetary policy and macroeconomic stabilisation.

This paper investigates whether the Phillips curve relationship carries forecasting value for Danish inflation. Denmark presents an interesting case because, as a small open economy operating under ERM II\footnote{ERM II links Denmark's krone to the euro at 7.46038 DKK/euro, aligning monetary policy with the ECB.}, Danish monetary conditions are closely tied to those of the euro area. This raises the question of whether broader European macroeconomic conditions carry additional forecasting value for Danish inflation beyond what domestic labour market variables can provide.

We address these questions through a sequential modelling strategy. First, we establish a univariate ARIMA benchmark estimated on the full training sample, 
and complement it with a second ARIMA model whose parameters are estimated 
exclusively on pre-inflation-surge data and subsequently conditioned on realised 
observations through the forecast origin. This pair of baselines allows us to 
assess whether the post-2021 inflation episode has altered the underlying 
time-series dynamics of Danish inflation. To understand the prediction impact of the external regressors, we will conduct Granger causality tests to examine 
whether EU27 inflation and Danish unemployment contain predictive information for Danish inflation beyond its own history. To implement the findings from the Granger causality test, we estimate a dynamic regression model that includes Danish unemployment directly as an exogenous regressor, providing a formal test of the domestic Phillips curve in the context of a forecasting framework. The results motivate a fourth model: a General-to-Specific (GETS) dynamic regression in which the General Unrestricted Model (GUM) encompasses domestic unemployment and EU27 inflation across multiple lags, with systematic selection determining which variables carry statistically meaningful predictive content. A surge dummy with bounds informed by a Quandt Likelihood Ratio (QLR) test is included as a selectable regressor in the GETS stage, allowing the data to determine whether a structural break correction is warranted, mirroring the motivation behind the pre-surge ARIMA model.

We used generative AI \parencite{anthropic2026} for idea generation and conceptualisation, such as brainstorming model framings and structuring the analysis.

The remainder of the paper is structured as follows. Section \ref{sec:literature_review} reviews the relevant literature. Section \ref{sec:data} describes the data. Section \ref{sec:methodology} introduces the methodology. Section \ref{sec:results} presents and discusses the results. Section \ref{sec:conclusion} concludes.